% !PRINTED 413
\weberpage{413}
If we here take $\alpha=\alpha_1$ and $\alpha=\alpha_2$, and set for
brevity $\varphi_{\alpha_1}=\varphi_1$, $\varphi_{\alpha_2}=\varphi_2$, then
from the two equations
\[
\varphi_1 = a_1\varphi + b_1\psi, \quad \varphi_2 = a_2\varphi + b_2\psi
\]
we can eliminate $\psi$ and obtain an equation of the form:
\weq{7}{h\varphi + h_1\varphi_1 + h_2\varphi_2 = 0,}
in which $h$, $h_1$, $h_2$ are constants, of which at least two are
different from zero. To the identical equation (7) we can now apply the
permutations $\alpha_1$, $\alpha_2$, and obtain, since
$\alpha_1\alpha_2 = \alpha_2\alpha_1 = \alpha_3$, $\alpha_1\alpha_1 = 1$,
$\alpha_2\alpha_2 = 1$,
\[
\begin{array}{l|l}
h\varphi + h_1\varphi_1 + h_2\varphi_2 = 0 & h,\\
h\varphi_1 + h_1\varphi + h_2\varphi_3 = 0 & h_1,\\
h\varphi_2 + h_1\varphi_3 + h_2\varphi = 0 & -h_2,
\end{array}
\]
from which, by multiplication with the factors standing beside them and
addition:
\weq{8}{(h^2 + h_1^2 - h_2^2)\varphi + 2hh_1\varphi_1 = 0.}
If, then, $h$ and $h_1$ are different from zero, $\varphi$ and $\varphi_1$
differ only by a constant factor (which must be $=\pm 1$, since $\alpha_1$
is of degree 2). But if $h$ or $h_1=0$, then from (7) it follows that
$\varphi_1=\pm\varphi_2$ or $\varphi=\pm\varphi_2$, and the application of
$\alpha_1$ to the first equation gives $\varphi=\pm\varphi_3$. Hence in
every case \gesperrt{one of the relations} holds:
\weq{9}{\varphi = \pm\varphi_1, \quad \varphi = \pm\varphi_2, \quad \varphi = \pm\varphi_3.}
Since now one can equally well regard $\gamma,\alpha_1$ as $\gamma,\alpha_2$,
as also $\gamma,\alpha_3$, as generating elements of the group $A$, this
result, taken together with (4), yields the theorem:

1. \gesperrt{The function $\varphi$ has the property of changing, under all
permutations of the alternating group, only by a constant factor.}

But the very same conclusion is applicable to $\psi$ as well, and hence
also to the quotient $\vartheta$ of the two functions. Accordingly, for
every permutation $\pi$ of the alternating group,
\weq{10}{\vartheta_\pi = \varepsilon\vartheta.}
Herein, when $\pi$ belongs to the cyclic permutations $\gamma$,
$\varepsilon$ is a third root of unity. But from this it is further to be
concluded, because all permutations $\pi$ are composed of $\gamma$, $\gamma'$
%CONT>>
